\chapter{Vergleich und Diskussion der Ergebnisse}
\label{chap:interpretation}

\section{Beantwortung der Forschungsfrage}

Die Forschungsfrage dieser Arbeit lautet: \enquote{Aus welchen Gründen verfolgen staatliche Akteure divergierende Strategien im Umgang mit Kryptowährungen?} Auf Grundlage der drei Fallstudien kann diese Frage weitgehend beantwortet werden. Die Analyse zeigt, dass die Unterschiede zwischen El Salvador, den Vereinigten Arabischen Emiraten und der Europäischen Union nicht zufällig entstehen, sondern aus unterschiedlichen politischen, wirtschaftlichen und institutionellen Ausgangslagen hervorgehen.

El Salvador verfolgte einen besonders offensiven Ansatz. Bitcoin wurde als gesetzliches Zahlungsmittel eingeführt, weil die Regierung damit finanzielle Inklusion, günstigere Rücküberweisungen und internationale Aufmerksamkeit erreichen wollte. Die Strategie beruhte stark auf der Hoffnung, dass eine neue Technologie bestehende strukturelle Probleme im Finanzsystem schnell lösen könne. Die Ergebnisse zeigen jedoch, dass diese Erwartung nur begrenzt erfüllt wurde. Die Nutzung blieb nach dem anfänglichen Interesse niedrig, viele Bürgerinnen und Bürger verwendeten die Chivo-App vor allem wegen des Startbonus, und Bitcoin konnte weder den US-Dollar ersetzen noch Rücküberweisungen in relevantem Umfang verändern.

Die VAE wählten dagegen keinen Ansatz, bei dem Kryptowährungen als normales Alltagsgeld eingeführt werden. Sie nutzen Krypto- und Blockchain-Technologien vor allem als Teil einer wirtschaftlichen Modernisierungs- und Standortstrategie. Die Regulierung durch VARA zeigt, dass die VAE digitale Vermögenswerte nicht grundsätzlich ablehnen, sie aber in einen kontrollierten Rechtsrahmen einbetten wollen. Das Ziel ist weniger die Ablösung bestehender Währungen, sondern die Stärkung des Finanz- und Technologiestandorts.

Die Europäische Union verfolgt wiederum einen vorsichtigen Regulierungsansatz. Mit MiCA werden Krypto-Assets rechtlich eingeordnet, beaufsichtigt und in das Finanzsystem eingebunden. Gleichzeitig bleibt deutlich, dass die EU private Kryptowährungen nicht als Ersatz für den Euro versteht. Der Schutz von Geldpolitik, Finanzstabilität, Verbraucherschutz und Vertrauen in öffentliches Geld steht im Vordergrund. Der geplante digitale Euro zeigt, dass die EU digitale Zahlungsformen nicht ablehnt, sondern eine staatlich kontrollierte Alternative bevorzugt.

Damit lässt sich die Forschungsfrage so beantworten: Staaten wählen unterschiedliche Strategien, weil sie Kryptowährungen durch die Brille ihrer eigenen Probleme und Interessen bewerten. Für El Salvador standen finanzielle Inklusion und internationale Sichtbarkeit im Mittelpunkt. Für die VAE waren Standortpolitik, Innovation und kontrollierte Marktöffnung entscheidend. Für die EU standen Regulierung, Verbraucherschutz und monetäre Stabilität im Vordergrund. Gemeinsam ist allen drei Fällen, dass Kryptowährungen nur dort akzeptiert werden, wo der Staat seine zentrale Steuerungsrolle nicht vollständig verliert.

\section{Einordnung in den Forschungsstand}

Die Ergebnisse bestätigen wesentliche Aussagen des bisherigen Forschungsstands. Die Forschung zu Bitcoin als Geld zeigt, dass Bitcoin häufig eher als Anlage- oder Spekulationsobjekt genutzt wird und nicht als alltägliches Zahlungsmittel. Genau dieses Muster wird besonders im Fall El Salvador sichtbar. Obwohl Bitcoin dort gesetzlich anerkannt wurde, entstand keine stabile Alltagsnutzung. Auch im europäischen Fall zeigt sich, dass Krypto-Assets zwar bekannt sind und teilweise gehalten werden, aber im Zahlungsverkehr nur eine begrenzte Rolle spielen.

Auch die Forschung zur Regulierung wird durch die Ergebnisse bestätigt. Staatliche Akteure reagieren nicht nur technisch oder wirtschaftlich auf Kryptowährungen, sondern vor allem institutionell. Sie fragen, wie Geldwäsche verhindert, Verbraucher geschützt, Anbieter beaufsichtigt und Finanzmärkte stabil gehalten werden können. Die VAE und die EU zeigen dabei zwei unterschiedliche, aber ähnliche Grundlogiken: Beide akzeptieren Krypto-Innovation nur unter Bedingungen staatlicher Kontrolle. Die konkrete Ausgestaltung ist verschieden, aber die Richtung ist vergleichbar.

Ebenso bestätigt die Analyse die Bedeutung von Vertrauen. Vertrauen entsteht nicht automatisch dadurch, dass eine Regierung eine Kryptowährung erlaubt oder sogar gesetzlich anerkennt. Im Fall El Salvador blieb das Vertrauen vieler Nutzer gering, weil Bitcoin stark schwankt, die App Probleme hatte und der Nutzen im Alltag begrenzt blieb. Im europäischen Fall zeigt sich ein ähnliches Grundproblem: Für alltägliche Zahlungen erwarten Menschen Stabilität, einfache Nutzung, rechtliche Sicherheit und breite Akzeptanz. Kryptowährungen erfüllen diese Erwartungen bisher nur teilweise.

Ein gewisser Widerspruch entsteht jedoch gegenüber optimistischen Annahmen, nach denen Kryptowährungen vor allem wegen ihrer technischen Vorteile schnell breite Akzeptanz finden könnten. Die untersuchten Fälle zeigen, dass technische Möglichkeiten allein nicht ausreichen. Selbst kostenlose Transaktionen, staatliche Förderung oder regulatorische Offenheit führen nicht automatisch zu Massennutzung. Entscheidend ist nicht nur, ob eine Technologie verfügbar ist, sondern ob sie im Alltag einen klaren Vorteil gegenüber bestehenden Zahlungsmitteln bietet.

Ein weiterer Widerspruch betrifft die Rolle des Staates. Kryptowährungen werden oft als dezentrale Alternative zu staatlichem Geld beschrieben. Die Fallstudien zeigen jedoch, dass ihre praktische Verbreitung stark von staatlichen Entscheidungen abhängt. El Salvador versuchte die Nutzung direkt durch Gesetz und App-Infrastruktur zu fördern. Die VAE und die EU schaffen regulatorische Räume, in denen Krypto-Anbieter tätig sein dürfen. Damit wird deutlich: Auch eine dezentrale Technologie bleibt in ihrer gesellschaftlichen Nutzung stark von politischen und rechtlichen Rahmenbedingungen abhängig.

\section{Weitere Erkenntnisprobleme}

Im Verlauf der Analyse wurden mehrere Erkenntnisprobleme deutlich. Erstens ist die Datenlage zur tatsächlichen Nutzung von Kryptowährungen uneinheitlich. Viele verfügbare Daten messen Downloads, Besitz oder registrierte Nutzer, aber nicht immer die dauerhafte Alltagsnutzung. Im Fall El Salvador ist zum Beispiel wichtig zu unterscheiden, ob Menschen die Chivo-App nur wegen des Bonus heruntergeladen haben oder ob sie Bitcoin danach regelmäßig verwendet haben.

Zweitens ist Vertrauen schwer messbar. In den Fallstudien zeigt sich zwar, dass fehlendes Vertrauen eine wichtige Rolle spielt, aber Vertrauen besteht aus mehreren Faktoren: technisches Verständnis, politische Glaubwürdigkeit, Sicherheit der Plattform, Preisstabilität und persönliche Erfahrung. Diese Faktoren lassen sich nicht allein durch Gesetzestexte oder allgemeine Nutzungszahlen erfassen.

Drittens unterscheiden sich die drei Untersuchungsräume stark in Größe, politischem System, wirtschaftlicher Struktur und Währungssituation. El Salvador ist ein dollarisiertes Land mit hoher Bedeutung von Rücküberweisungen. Die VAE sind ein finanzstarker Standortstaat mit globaler Investitionsstrategie. Die EU ist ein großer Währungsraum mit gemeinsamer Zentralbank und komplexer Regulierung. Der Vergleich ist deshalb erkenntnisreich, aber nicht vollständig symmetrisch.

Viertens bleibt offen, wie sich die Strategien langfristig entwickeln. MiCA ist noch relativ neu, der digitale Euro befindet sich weiterhin in Vorbereitung, und auch der Krypto-Standort Dubai entwickelt sich weiter. Deshalb können die Ergebnisse vor allem den bisherigen Stand erklären, aber nur begrenzt vorhersagen, welche Strategie langfristig erfolgreicher sein wird.

\section{Empfehlungen für weitere Forschung}

Für weitere Forschung wäre zunächst eine stärkere empirische Untersuchung der tatsächlichen Nutzung sinnvoll. Besonders hilfreich wären Umfragen und Interviews mit Bürgerinnen und Bürgern, Händlern und Unternehmen. Dadurch ließe sich genauer erfassen, warum Menschen Kryptowährungen nutzen, warum sie sie ablehnen und welche Rolle Vertrauen, Bildung, Einkommen und technische Erfahrung spielen.

Zweitens wären vergleichende Langzeitstudien sinnvoll. Sie könnten untersuchen, ob Krypto-Regulierungen wie MiCA oder VARA langfristig zu mehr Sicherheit, mehr Innovation oder stärkerer Nutzung führen. Dabei sollte nicht nur die Zahl der Anbieter betrachtet werden, sondern auch die tatsächliche Nutzung im Zahlungsverkehr und die Auswirkungen auf Verbraucher.

Drittens sollte die Forschung stärker zwischen verschiedenen Formen digitaler Währungen unterscheiden. Private Kryptowährungen, Stablecoins und digitales Zentralbankgeld erfüllen unterschiedliche Funktionen und werden vom Staat unterschiedlich bewertet. Besonders der digitale Euro könnte künftig zeigen, ob Menschen digitalen Zahlungsmitteln eher vertrauen, wenn sie öffentlich kontrolliert und wertstabil sind.

Viertens wäre eine stärkere Verbindung quantitativer und qualitativer Methoden sinnvoll. Nutzungsstatistiken können zeigen, wie stark Kryptowährungen verwendet werden. Interviews und Fallstudien können erklären, warum diese Nutzung entsteht oder ausbleibt. Eine Kombination beider Methoden wäre geeignet, um die Lücke zwischen technischer Möglichkeit, staatlicher Regulierung und gesellschaftlichem Vertrauen genauer zu verstehen.

Insgesamt zeigt die Interpretation, dass Kryptowährungen bisher nicht an einer einzelnen Ursache scheitern. Vielmehr treffen technische Risiken, staatliche Kontrollinteressen und begrenztes gesellschaftliches Vertrauen zusammen. Gerade deshalb bleiben sie in den untersuchten Fällen eher regulierte Vermögenswerte oder politische Innovationsprojekte als normales Geld des Alltags.
